% No preamble. Prefix mrg-. Homogeneous field domains, not a marked RG theorem.
\section{Nonlinear current control and a finite-block radius gain}
\label{mrg-section}
Fix the primitive odd ratio $\ell$, gauge group and $0<\alpha<1$.
All graphs are finite compatible periodic regulators. Constants are
independent of preceding refinement $M=\ell^k$ and volume; constants
subscripted $q$ may depend on a fixed number $q$ of primitive steps.
All complex functions use the branches through the identity.

\begin{proposition}[Physical current of the homogeneous minimum]
\label{mrg-current}
Use Hermitian Lie coordinates, normalized trace $\tau$, physical pairing
$\langle v,w\rangle_\eta=\eta^d\sum_e\tau(v_ew_e)$, and
\[
 S_\eta(U)=\eta^{d-4}\sum_p\{1-\tau(P_p+P_p^{-1})/2\},\qquad
 \left.\frac d{dt}S_\eta(e^{i\eta tv}U)\right|_{t=0}
       =\langle J_{\rm phys}(U),v\rangle_\eta .
\]
The sum counts each unoriented plaquette once. Explicitly,
\begin{equation}
 J_{\rm phys}(U)(e)=\eta^{-3}\sum_{p\ni e}
       \pi_{\mathfrak g}\frac{P_{p,e}-P_{p,e}^{-1}}{2i},
 \qquad
 \|J_{\rm phys}(U_{\rm ord}(B))\|_\infty\le C\|B\|_\infty.
 \label{mrg-current-bound}
\end{equation}
Here $P_{p,e}$ is based and oriented with first edge the positive bond
$e$, and repeated incidences count with multiplicity. The last map is
holomorphic on a fixed coarse product ball, uniformly in $M$ and volume.
\end{proposition}
\begin{proof}
Write the corrected constraint as $\mathcal F(A)=QA+\mathcal R(A)$,
its minimum as $A_*(B)$, and its potential Euler derivative as
$E(A)=KA+N(A)$. Notebook15's full tangent decomposition
\texttt{ng-tangent-decomposition} and right inverse of $D\mathcal F$
give $E(A_*)=D\mathcal F(A_*)^Tm$. Applying $H^T$ identifies
the multiplier exactly as \texttt{nm-multiplier}:
\[
 m=(I+H^TD\mathcal R(A_*)^T)^{-1}
       \{q[B-\mathcal R(A_*)]+H^TN(A_*)\}.
\]
The bounds \texttt{nm-adjoint-composite} and
\texttt{nm-multiplier-bounds} give $\|E(A_*)\|_\infty\le C\|B\|$.
The real identity extends holomorphically to the common complex ball.
If $\partial_t e^{i(X+tY)}|_0=iT_X(Y)e^{iX}$, then
$E(A)=T_{\eta A}^TJ_{\rm phys}(e^{i\eta A})$.
The inverse transpose of $T_{\eta A}$ is uniformly bounded for small
$A$ and $\eta\le1$. Finally \texttt{ot-ordinary-representative}
conjugates the current by a bounded holomorphic gauge and its inverse.
This proves the bound without estimating a rough outside divergence.
Differentiating the displayed action gives the stated normalization.
\end{proof}

\paragraph{Explicit source-domain hypothesis (UF).}
For the family of radii $(\alpha_0,\alpha_1)$ being used, assume that
source conditions (i)--(iii), with both radii reduced by one fixed
factor $c_{\rm src}>0$, imply the reconstructed-domain condition (iv),
uniformly as the radii shrink and for the homogeneous restrictions in
question. The smaller-constant paragraph of CMP109, pp.~262--263,
alone is a fixed-radii statement and does not establish this uniformity.

\begin{corollary}[Conditional source specialization]
\label{mrg-specialization}
Under UF, if $R_X(U,J)$ is holomorphic and bounded by $A_X$ on that
source domain, then
\begin{equation}
 f_X(B)=R_X(U_{\rm ord}(B),J_{\rm phys}(U_{\rm ord}(B))),\qquad
 |f_X(B)|\le A_X\quad
 (\|B\|_\infty<c\min\{\alpha_0,\alpha_1,b_*\}).
 \label{mrg-source-ball}
\end{equation}
The constant $c>0$ is independent of $M,X$ and volume.
\end{corollary}
\begin{proof}
The ordinary potential has physical $C^{1,\alpha}$ norm at most
$C\|B\|$ by \texttt{ot-ordinary-bound}. Plaquette multiplication gives
$\|U(\partial p)-I\|\le C\eta^2(\|B\|+\|B\|^2)$, and
Proposition~\ref{mrg-current} controls the current. Use real reference
$I$ for the complex perturbation, choose $c$ to fit the reduced radii,
and apply UF. Every collar is restricted from the same global field.
\end{proof}
This composition is generally nonlocal in $B$. It does not provide the
source's independent-current or decoupling-parameter localization norm.

\begin{proposition}[Exact noisy charts and fixed-block contraction]
\label{mrg-block}
Let $P_{M,R}(e^{iw})$ be the actual ordinary prefix of the $MR$-refined
minimum and $a_{M,R}=\log P_{M,R}/i$. The prefix theorem
\texttt{pgh-scale-bounds} gives $\|a_{M,R}(w)\|\le C_0\|w\|/R$,
with $C_0$ independent of $R=\ell^s$.
Root $V=P_{M,\ell}(e^{iw})$ by its block contour gauge $h(w)$,
where $h(y)=I$ at coarse centers. For the actual retained pivot chart
of Notebook19 define
\begin{equation}
 F_M(w,t)=\frac1i\log\left[
       \{e^{i\Phi_{V^{h(w)}}(t)}V^{h(w)}\}^{h(w)^{-1}}\right].
 \label{mrg-chart}
\end{equation}
On fixed product balls this is holomorphic and
\begin{equation}
 F_M(w,0)=a_{M,\ell}(w),\quad
 \|F_M(w,t)-F_M(w,0)\|\le C_\ell\|t\|,\quad
 \|D_wF_M(w,t)\|_{\infty\to\infty}\le C_{\rm prim}.
 \label{mrg-primitive}
\end{equation}
For each fixed $q$, compose these exact maps, with depths
$M,M\ell,\ldots,M\ell^{q-1}$, to obtain $\mathcal B_{M,q}(w;t)$.
There are fixed $b_q,\rho_q>0$ and a bounded holomorphic gauge
$\gamma_{M,q}(w)$, independent of all $t_a$, such that
\begin{equation}
 \widetilde{\mathcal B}_{M,q}
   =\frac1i\log[(e^{i\mathcal B_{M,q}})^{\gamma_{M,q}(w)}],
 \qquad
 \|\widetilde{\mathcal B}_{M,q}(w;t)\|
       \le C_0\ell^{-q}\|w\|+C_q\max_{a<q}\|t_a\|
 \label{mrg-block-gain}
\end{equation}
on $\|w\|<b_q$, $\max_a\|t_a\|<\rho_q$.
\end{proposition}
\begin{proof}
Undoing the same background gauge gives the exact link expression
$e^{i\operatorname{Ad}_{h(e_-)^{-1}}\Phi(t)}e^{ia_{M,\ell}(w)}$.
The retained-chart bound $\|\Phi(t)\|\le C_\ell\|t\|$ and the
uniform logarithm derivative prove the first two claims of
\eqref{mrg-primitive}; Banach Cauchy on a fixed larger ball proves
the third. No coordinate count enters.
For fixed $q$, shrink the input and noise balls to keep every
intermediate field in these balls. Telescoping gives
\[
 \|\mathcal B_{M,q}(w;t)-\mathcal B_{M,q}(w;0)\|
       \le C_\ell\sum_{a<q}C_{\rm prim}^{a}\|t_a\|.
\]
Notebook31's two-minimum comparison \texttt{csg-comparisons}
identifies the zero-noise composition modulo gauge with
$P_{M,\ell^q}(e^{iw})$. A holomorphic gauge choice follows from its
explicit construction: ordinary-to-corrected transfer, flow into the
fixed slice, and inverse transfer. The transfers are holomorphic.
The complexified flow \texttt{ng-flow} has the exact gauge-error
identity $R_0D^TA(t)=e^{-t}R_0D^TA(0)$ and the uniform bound
$\|A'(t)\|_{C^{1,\alpha}}\le Ce^{-t}\|R_0D^TA(0)\|_{C^\alpha}$.
On a smaller complex ball this gives global continuation and uniform
holomorphic convergence. The scalar gauge generator has the same
integrable $C^{2,\alpha}$ bound, so its gauge and inverse converge
holomorphically as well. The finitely many scale conversions in
\texttt{csg-rescaling} cost only $q$-dependent constants. Their
composition supplies $\gamma_{M,q}(w)=I+O_q(\|w\|)$ uniformly in
$M$ and volume. The real minimum identity extends to the complex ball.
Applying this same gauge to the noisy composition changes its difference
from zero noise by a bounded factor. The direct prefix bound proves
\eqref{mrg-block-gain}.
\end{proof}

Choose $q$ once with $\ell^q>2C_0$ and put $\theta=C_0\ell^{-q}<1/2$.
For actual retained variables $t_a=g_{j+a}z_a$ the noise cost is
$e_j=C_q\max_{a<q}g_{j+a}R_{\rm cut}(g_{j+a})$.
Thus, for $1<\lambda<\theta^{-1}$, the block-endpoint recursion
\begin{equation}
 r_{n+1}=\min\{\lambda r_n,b_q\},\qquad
 e_{j+nq}\le\tfrac12(1-\theta\lambda)r_n
 \label{mrg-radius-recursion}
\end{equation}
has a fixed fractional domain margin and reaches $b_q$ in
$O(1+\log(b_q/r_0))$ blocks. If $r_0=cg_i(\log g_i^{-2})^s$,
$R_{\rm cut}(g)\le C(\log g^{-2})^p$ with $s>p$, and
$|g_{j+1}^{-2}-g_j^{-2}|\le D$ on the bootstrap prefix, its noise
condition holds for small $g_i$: over $O(\log(1/g_i))$ steps,
$g_j\asymp g_i$, while $e_j/r_0=O((\log g_i^{-2})^{p-s})$.
For $s=p$ the available constants must be checked; $s<p$ gives no
startup by this argument. Source parameter order and UF remain inputs.

\paragraph{Exact kernel ledger: an additional matching condition.}
For a particular retained source kernel require, for every invariant
scalar $f$, the exact identity
\[
 (K_jf)(w)=\int a_j(w,z)f(F_{M_j}(w,g_jz))\,dz,
 \qquad K_jf\ \hbox{gauge invariant}.
\]
Here $a_j$ retains every rooted weight, Jacobian, covariance, cutoff
and normalization. Unrooting changes only the argument of $f$.
Applying this identity to the invariant scalar $K_{j-1}f$ and then
substituting its integral yields the composed $F$ maps of
Proposition~\ref{mrg-block}; earlier weights are evaluated at the new
coarse argument, and no gauge is commuted through them. This elementary
iteration, with the requisite finite-cutoff Fubini majorants, proves
the ledger if the displayed conditions hold. Covariance of a full
conditional does not alone imply covariance of an arbitrary retained
restriction. For the original regular source kernel, CMP109, p.~269,
(2.16) and the adjacent paragraph establish covariance of the whole
measure and cutoff by local orthogonal conjugations of the retained
field. An invariant scalar insertion preserves that identity, as does
a scalar radial cutoff on the same retained bonds. Other restrictions
require their own check; no mixed-sector matching is inferred.

\paragraph{Scope of the gain.}
Proposition~\ref{mrg-block} proves the composed homogeneous field-domain
bound independently of that kernel match. Its use for a scalar requires
gauge invariance on a compatible gauge-orbit domain. No single
large-ratio quantum step replaces the separate kernels.
With fractional margin, Banach Cauchy gives
$|f(B(w,z))-f(B(w,0))|\le C A_XgR_{\rm cut}/r$; converting this
into a localized prepared-potential norm still requires the source
support and independent $(U,J)$ decoupling bounds. The uniform marked
history estimate, including regenerated marginal coefficients and their
coupling feedback, is not proved by this radius gain.
